\documentclass[../../main.tex]{subfiles}% 注意这里的文件路径不能用 ./main.tex ,否则用latexmk编译子文件会报错
\graphicspath{{\subfix{./image/}}{\subfix{../../image/}}} % 指定图片引用路径,后续可以直接使用图片文件名
% 如果图片存放在上级目录,则使用 ../../image/ 作为备用路径

% 例如:
% \begin{figure}[H]
% \centering
% \includegraphics[width=0.3\textwidth]{图.png}
% \caption{}
% \label{figure:图}
% \end{figure}
% 注意:上述\label{}一定要放在\caption{}之后,否则引用图片序号会只会显示??.

\begin{document}

\section{正测度集与矩体的关系}

\begin{theorem}\label{theorem:定理2.19}
设\(E\)是\(\mathbb{R}^n\)中的可测集，且\(m(E)>0\)，\(0 < \lambda < 1\)，则存在开矩体\(I\)，使得
\begin{align}
m(I\cap E)>\lambda|I|.\label{eq::::--FSAFF-2.8}
\end{align}
\end{theorem}
\begin{remark}
上述定理告诉我们，任何一个正测集，其中总有一部分被一个矩体套住，使两者的测度差小于预先给定的正数\(\varepsilon\)。当然，这一测度差不一定能等于零。 
\end{remark}
\begin{proof}

{\heiti 情形I:}当\(m(E)<+\infty\)时,对于\(0 < \varepsilon < (\lambda^{-1}-1)m(E)\)，作\(E\)的\(L\)-覆盖\(\{I_k\}\)，使得
\begin{align*}
\sum_{k = 1}^{\infty}|I_k|<m(E)+\varepsilon.
\end{align*}
从而存在\(k_0\)，使得\(\lambda|I_{k_0}|<m(I_{k_0}\cap E)\)。事实上，若对一切\(k\)，有
\begin{align*}
\lambda|I_k|\geqslant m(I_k\cap E),
\end{align*}
则可得
\begin{align*}
m(E)=m(E\cap \bigcup_{k=1}^{\infty}I_k)\leqslant\sum_{k = 1}^{\infty}m(I_k\cap E)\leqslant\lambda\sum_{k = 1}^{\infty}|I_k|\leqslant\lambda(m(E)+\varepsilon)<m(E).
\end{align*}
这就导致\(m(E)<m(E)\)，产生矛盾!

{\heiti 情形II:}当$m(E)=+\infty$时,由\nrefthe{theorem:定理2.13}{(ii)}可知，存在闭集$F\subset E$,使得$m(E\backslash F)<1$.记$H=E\backslash F$,则$m(H)<1$且$H\subset E$.于是由{\heiti 情形I}可知,存在开矩体$I$,使得
\begin{align*}
\lambda |I|<m(I\cap H).
\end{align*}
再由$I\cap H\subset I\cap E$及测度的单调性可得
\begin{align*}
\lambda |I|<m(I\cap H)\leqslant m(I\cap E).
\end{align*}
故结论得证.
\end{proof}

\begin{proposition}
设$E_1,E_2\subseteq \mathbb{R}^n$是正测度集（$m(E_1)>0,m(E_2)>0$）。则存在向量$h\in \mathbb{R}^n$使得
\begin{align*}
m\bigl(E_1\cap (E_2+h)\bigr) > 0.
\end{align*}
\end{proposition}
\begin{proof}
反证. 假设对$\forall h\in \mathbb{R}^n$，有
\begin{align}\label{eq_311_9x2z}
m\left( E_1\cap \left( E_2+h \right) \right) =0.
\end{align}
由\refthe{theorem:定理2.19}知对$\forall \lambda \in \left( 0,1 \right)$，存在开矩体$I_1,I_2$，使得
\begin{align}\label{eq_312_5b7m}
m\left( E_1\cap I_1 \right) >\lambda \left| I_1 \right|,\quad m\left( E_2\cap I_2 \right) >\lambda \left| I_2 \right|.
\end{align}
设$I_1,I_2$的中心分别为$c_1,c_2$，$I_1,I_2$的各边长分别为$a_1,a_2,\cdots ,a_n$和$b_1,b_2,\cdots ,b_n$. 取
$$h=c_1-c_2,\quad I_{2}^{\prime}=I_2+h,\quad r=\frac{\min\limits_{i=1,2,\cdots ,n}\{a_i,b_i\}}{2},$$
则$I_{2}^{\prime}$与$I_1$的中心重合，且$I_1\cap I_{2}^{\prime}\supseteq B\left( c_1,r \right)$. 由\hyperref[theorem:外测度的平移不变性]{测度的平移不变性}知
$$\left| I_2 \right|=m\left( I_2 \right) =\left| I_{2}^{\prime} \right|,\quad m\left( E_2\cap I_{2}^{\prime} \right) =m\left( \left( E_2\cap I_2 \right) +h \right) =m\left( E_2\cap I_2 \right).$$
再结合\rrefthe{theorem:测度的基本性质}{theorem:测度的基本性质-4}知
\begin{align}\label{eq_313_4d8s}
m\left( \left( E_1\cap I_1 \right) \cup \left( E_2\cap I_{2}^{\prime} \right) \right)
\leqslant m\left( I_1\cup I_2' \right)
= m\left( I_1 \right) +m\left( I_{2}^{\prime} \right) -m\left( I_1\cap I_{2}^{\prime} \right)
\leqslant \left| I_1 \right|+\left| I_2 \right|-m\left( B\left( c_1,r \right) \right).
\end{align}
又由\rrefthe{theorem:测度的基本性质}{theorem:测度的基本性质-4}和\eqref{eq_311_9x2z}\eqref{eq_312_5b7m}式可得
\begin{align}\label{eq_314_6f9g}
m\left( \left( E_1\cap I_1 \right) \cup \left( E_2\cap I_{2}^{\prime} \right) \right)
&= m\left( E_1\cap I_1 \right) +m\left( E_2\cap I_{2}^{\prime} \right) -m\left( \left( E_1\cap I_1 \right) \cap \left( E_2\cap I_{2}^{\prime} \right) \right) \nonumber\\
&\geqslant m\left( E_1\cap I_1 \right) +m\left( E_2\cap I_2 \right) -m\left( E_1\cap \left( \left( E_2\cap I_2 \right) +h \right) \right) \nonumber\\
&= m\left( E_1\cap I_1 \right) +m\left( E_2\cap I_2 \right) -m\left( E_1\cap \left( E_2+h \right) \right) \nonumber\\
&= m\left( E_1\cap I_1 \right) +m\left( E_2\cap I_2 \right)
\geqslant \lambda \left( \left| I_1 \right|+\left| I_2 \right| \right).
\end{align}
取$\lambda =1-\frac{m\left( B\left( c_1,r \right) \right)}{2\left( \left| I_1 \right|+\left| I_2 \right| \right)}$，则由\eqref{eq_313_4d8s}\eqref{eq_314_6f9g}式知
$$\lambda \left( \left| I_1 \right|+\left| I_2 \right| \right) \leqslant \left| I_1 \right|+\left| I_2 \right|-m\left( B\left( c_1,r \right) \right)
\Longrightarrow 1-\frac{m\left( B\left( c_1,r \right) \right)}{2\left( \left| I_1 \right|+\left| I_2 \right| \right)}=\lambda \leqslant 1-\frac{m\left( B\left( c_1,r \right) \right)}{\left( \left| I_1 \right|+\left| I_2 \right| \right)},$$
显然矛盾!
\end{proof}


\begin{proposition}\label{proposition:正测集与区间交测度小于区间}
\([0,1]\)中存在正测集\(E\)，使对\([0,1]\)中任一开区间\(I\)，有
\begin{align*}
0 < m(E\cap I) < m(I).
\end{align*}
\end{proposition}
\begin{remark}
证明中的邻接区间是指构造类Cantor集时挖去的区间.
\end{remark}
\begin{proof}
首先，在\([0,1]\)中作类Cantor集\(H_1\)：\(m(H_1)=\frac{1}{2}\)。其次，在\([0,1]\)中\(H_1\)的邻接区间\(\{I_{1j}\}\)的每个\(I_{1j}\)内再作类Cantor集\(H_{1j}\)：
\(m(H_{1j}) = \frac{|I_{1j}|}{2^2}\)，并记\(H_2 = \bigcup_{j = 1}^{\infty}H_{1j}\)。然后，对\(H_1\cup H_2\)的邻接区间\(\{I_{2j}\}\)的每个\(I_{2j}\)，又作类Cantor集\(H_{2j}\)：\(m(H_{2j}) =\frac{|I_{2j}|}{2^3}\)。再记\(H_3 = \bigcup_{j = 1}^{\infty}H_{2j}\)，依次继续进行，则得\(\{H_m\}\)。令\(E = \bigcup_{n = 1}^{\infty}H_n\)，得证。 
\end{proof}

\begin{definition}[向量差集]\label{definition:向量差集}
设$A,B$为两个非空集合,定义$A,B$的\textbf{向量差集}为
\begin{align*}
A - B \stackrel{\text{def}}{=} \{x - y: x\in A,y \in B\}.
\end{align*}
\end{definition}

\begin{example}
设 \(A, B\) 是 \(\mathbb{R}^n\) 的子集，定义 \(A+B = \{a+b \mid a \in A, b \in B\}\)。
\begin{enumerate}[(1)]
\item 证明：若 \(A\) 是开集，则 \(A+B\) 也是开集；
\item 给出一个例子，使得 \(A, B\) 都是闭集，但 \(A+B\) 不是闭集。
\end{enumerate}
\end{example}
\begin{proof}
\begin{enumerate}[(1)]
\item 设 \(z \in A+B\)。根据定义，存在 \(a \in A\) 和 \(b \in B\)，使得 \(z = a+b\)。因为 \(A\) 是开集，所以存在 \(r>0\)，使得开球 \(B(a, r) \subseteq A\)。考虑开球 \(B(z, r)\)。对于任意 \(y \in B(z, r)\)，我们有
\begin{align*}
\|y - z\| < r \implies \|(y - b) - a\| = \|y - (a+b)\| < r.
\end{align*}
因此 \(y - b \in B(a, r) \subseteq A\)，从而 \(y = (y-b) + b \in A+B\)。这说明 \(B(z, r) \subseteq A+B\)。由于 \(z\) 的任意性，\(A+B\) 是开集。

\item 取 \(A = \{(x, y) \in \mathbb{R}^2 \mid x>0, xy \geqslant 1\}\)，\(B = \{(x, y) \in \mathbb{R}^2 \mid x=0\}\)。

先证明 \(A\) 是闭集。若 \(A\) 中的点列 \((x_n, y_n) \to (x_0, y_0)\)，则 \(x_n>0\) 且 \(x_n y_n \geqslant 1\)。若 \(x_0=0\)，则 \(x_n \to 0\)，从而 \(y_n \to +\infty\)，这与 \((x_n, y_n)\) 收敛于有限点矛盾。故必有 \(x_0 > 0\)。对不等式两边取极限可得 \(x_0 y_0 \geqslant 1\)，即 \((x_0, y_0) \in A\)，所以 \(A\) 是闭集。
集合 \(B\) 显然是 \(y\) 轴，为闭集。

接着计算 \(A+B\)。由定义可知，\(A+B = \{(x, y) \in \mathbb{R}^2 \mid x>0\}\)。该集合是右半开平面，它是一个开集.因为点列 \((\frac{1}{n}, 0) \in A+B\) 当 \(n \to \infty\) 时收敛于 \((0,0)\)，但 \((0,0) \notin A+B\)，因此 \(A+B\) 不是闭集。
\end{enumerate}
\end{proof}

\begin{theorem}[Steinhaus定理]\label{theorem:Steinhaus定理}
设\(E\)是\(\mathbb{R}^n\)中的可测集，且\(m(E)>0\)。作（向量差）点集
\begin{align*}
E - E \stackrel{\text{def}}{=} \{x - y: x,y \in E\},
\end{align*}
则存在\(\delta_0>0\)，使得\(E - E \supset B(0,\delta_0)\)。
\end{theorem}
\begin{proof}
取\(\lambda\)满足\(1 - 2^{-(n + 1)}<\lambda<1\)$(n\geqslant  2)$。由\refthe{theorem:定理2.19}可知，存在矩体\(I\)，使得\(\lambda|I|<m(I\cap E)\)。现在记\(I\)的最短边长为\(\delta\)，并作开矩体
\begin{align*}
J = \left\{x = (\xi_1,\xi_2,\cdots,\xi_n): |\xi_i|<\frac{\delta}{2}\ (i = 1,2,\cdots,n)\right\}.
\end{align*}
从而只需证明\(J\subset E - E\)即可(在$J$中任取一个以原点为中心的开球$B(0,\delta_0)$)，也就是只要证明对每个\(x_0\in J\)，点集\(E\cap I\)必与点集\((E\cap I)+\{x_0\}\)相交（此时任取$y\in (E\cap I)\cap((E\cap I)+\{x_0\})$,从而存在$z\in E\cap I$,使得$y=z+x_0$.也即存在\(y,z\in E\cap I\subset E\)，使得\(y - z = x_0\)）。因为\(J\)是以原点为中心，边长为\(\delta\)的开矩阵，所以\(I\)的平移矩体\(I+\{x_0\}\)仍含有\(I\)的中心，从而知
\begin{align*}
m(I\cap (I+\{x_0\}))>2^{-n}|I|(n\geqslant  2).
\end{align*}
结合上式,再由\nrefthe{theorem:测度的基本性质}{(4)}可得
\begin{align*}
m(I\cup (I+\{x_0\}))=|I|+m(I+\{x_0\})-m(I\cap (I+\{x_0\}))
<2|I|-2^{-n}|I|,
\end{align*}
即
\begin{align*}
m(I\cup (I+\{x_0\}))<2\lambda|I|.
\end{align*}
但由于\(E\cap I\)与\((E\cap I)+\{x_0\}\)有着相同的测度并且都大于\(\lambda|I|\)，同时又都含于\(I\cup (I+\{x_0\})\)之中，故它们必定相交，否则其并集测度要大于\(2\lambda|I|\)，从而引起矛盾。
\end{proof} 

\begin{proposition}\label{proposition:存在某个正测集的Cauchy方程函数必是线性函数}
设有定义在 $\mathbb{R}$ 上的函数 $f(x)$, 满足
\begin{align*}
f(x + y) = f(x) + f(y),\quad x,y\in\mathbb{R},
\end{align*}
且在 $E\subset\mathbb{R}$ ($m(E)>0$) 上有界, 则 $f(x)=cx$ ($x\in\mathbb{R}$), 其中 $c = f(1)$.
\end{proposition}
\begin{proof}

（i）首先, 由题设知, 对 $r\in\mathbb{Q}$, 必有 $f(r)=rf(1)$.

（ii）其次, 由 $m(E)>0$ 可知, 存在区间 $I$: $I\subset E - E$. 不妨设 $|f(x)|\leqslant M$ ($x\in E$), 又对任意的 $x\in I$, 有 $x',x''\in E$, 使得 $x = x' - x''$, 则
\[
|f(x)| = |f(x') - f(x'')|\leqslant |f(x')| + |f(x'')|\leqslant 2M.
\]

记 $I = [a,b]$, 并考查 $[0,b - a]$. 若 $x\in [0,b - a]$, 则 $x + a\in [a,b]$. 从而由 $f(x) = f(x + a) - f(a)$ 可知, $|f(x)|\leqslant 4M$, $x\in [0,b - a]$. 记 $b - a = c$, 这说明
\begin{align*}
|f(x)|\leqslant 4M,\quad x\in [0,c].
\end{align*}
易知
\begin{align*}
|f(x)|\leqslant 4M,\quad x\in [-c,c].
\end{align*}

已知对任意的 $x\in\mathbb{R}$ 以及自然数 $n$, 均存在有理数 $r$, 使得 $|x - r|<\frac{c}{n}$,从而
\begin{align*}
n\left| f\left( x-r \right) \right|=\left| f\left( n\left( x-r \right) \right) \right|\leqslant 4M\Longrightarrow \left| f\left( x-r \right) \right|\leqslant \frac{4M}{n}.
\end{align*}
因此我们得到
\begin{align*}
|f(x)-xf(1)|&=\left| f\left( x-r \right) +f\left( r \right) -xf\left( 1 \right) \right|=\left| f\left( x-r \right) +rf\left( 1 \right) -xf\left( 1 \right) \right|
\\
&=\left| f\left( x-r \right) +\left( r-x \right) f\left( 1 \right) \right|\leqslant \left| f\left( x-r \right) \right|+\left| r-x \right|\left| f\left( 1 \right) \right|
\\
&\leqslant \frac{4M+c|f(1)|}{n}.
\end{align*}
根据 $n$ 的任意性 ($r$ 的任意性), 即得 $f(x)=xf(1)$. 
\end{proof}




















\end{document}